\begin{lemma}
\label{editorial-source-lemma-radicial-integral}
Let $\varphi : R \to S$ be a ring map. Assume
\begin{enumerate}
\item $\varphi$ is integral,
\item $\varphi$ induces an injective map of spectra,
\item $\varphi$ induces purely inseparable residue field extensions.
\end{enumerate}
Then $\varphi$ induces a homeomorphism from $\Spec(S)$ onto a closed
subset of $\Spec(R)$ and for any ring map
$R \to R'$ properties (1), (2), (3) are true for $R' \to R' \otimes_R S$.
\end{lemma}

\begin{proof}
Put $I=\ker\varphi$. The original integral
injection $R/I\to S$ gives image $V(I)$
on spectra by lying over. More generally,
for every ideal $J\subseteq S$, the
map $R/\varphi^{-1}(J)\to S/J$ is an
integral injection, since its monic
equations are the images of the
original ones. Its spectrum is
surjective, so the image of $V(J)$
is exactly $V(\varphi^{-1}(J))$.
Thus the original spectrum map is
closed. Its assumed injectivity
makes it a homeomorphism onto
the closed subset $V(I)$: closed
subsets of the domain have closed
images in this subspace, which
proves continuity of the inverse.

For any original $R\to R'$, write
$S'=R'\otimes_R S$. Each original
monic equation for $s\in S$ maps
to the same monic equation for
$1\otimes s$ with coefficients
mapped into $R'$. Those elements
generate $S'$ over $R'$, so
integrality is preserved, including
infinite generating families by
the finite-support subring argument.
Lemma \ref{editorial-source-lemma-radicial} supplies
injectivity on spectra and pure
inseparability of residue fields
for this exact base-changed map.
The preceding closed-image proof
therefore applies to it as well.

There is also an exact radical
comparison for its kernel $I'$:
$$
\sqrt{I'}=\sqrt{IR'}.
$$
Indeed the original image is
$V(I)$, and the explicit fibre
construction in Lemma
\ref{editorial-source-lemma-radicial} identifies the
new image with its inverse image
in $\Spec(R')$, namely $V(IR')$.
Integrality also identifies the
new image with $V(I')$. Equality
of these closed sets gives the
stated equality of radicals.
The same formula in fact holds
for every integral ring map:
for a prime of $R'$ containing
$IR'$, its contraction contains
$I$, so lying over supplies a
prime of $S$ above that contraction;
the nonzero tensor of the two
original residue fields supplies
a prime after base change. This
proves equality of the images
without requiring uniqueness.
Equality of the kernels themselves
need not hold: the explicit cusp
base change in Lemma
\ref{editorial-source-lemma-2-3-ring-map} has $I=0$
and $I'=(b)\neq0$ in $k[b]/(b^2)$,
while both radicals are $(b)$.
All comparisons concern the
original extended ideal $IR'$.
\end{proof}